\documentclass[10pt]{article}
\usepackage[margin=0.85in]{geometry}
\usepackage{amsmath,amssymb,amsthm,hyperref}
\newtheorem{theorem}{Theorem}
\newcommand{\SLP}{\mathcal{SL\text{-}P}^{+}}
\title{A counterexample to positive shifted Laguerre--P\'olya preservation}
\author{Principia Math in Progress}
\date{13 September 2026}
\begin{document}
\maketitle

\paragraph{Task 1. Resolve the Laguerre-operator preservation conjecture.}
Wagner's Conjecture 4.3 asserts that \(L_j(\phi)\in\SLP\) for every
\(\phi\in\SLP\) and every positive integer \(j\) \cite{W}.
The operators are defined by
\[
 |\phi(x+iy)|^2=\sum_{j\ge0}L_j(\phi)(x)y^{2j},\qquad
 L_1(\phi)=(\phi')^2-\phi\phi''.
\]
We disprove this assertion already for \(j=1\).

Recall the defining conditions in \cite[p.~3]{W}. Membership in the
shifted class of type I requires real polynomials \(p_N\), of degrees
\(d_N\), converging locally uniformly to \(\phi\), such that for
each fixed \(d\), the polynomial \(p_N^{(d_N-d)}\) has all its roots
real and of the same sign for every sufficiently large \(N\).
Membership in \(\SLP\) additionally requires the exponential Taylor
coefficients \(\gamma_n=\phi^{(n)}(0)\) to be eventually nonnegative.

\begin{theorem}
For any \(c>0\) and \(a>0\), the function \(\phi(x)=c+e^{ax}\)
belongs to \(\SLP\), whereas \(L_j(\phi)\notin\SLP\) for every
odd positive integer \(j\). In particular, \(\phi(x)=1+e^x\)
is a counterexample to Conjecture 4.3.
\end{theorem}

\begin{proof}
Set \(p_N(x)=c+(1+ax/N)^N\). These real polynomials have degree
\(N\) and converge to \(c+e^{ax}\) uniformly on each compact subset
of \(\mathbb C\). For fixed \(d\ge1\) and \(N>d\), differentiation
annihilates the constant term and gives
\[
 p_N^{(N-d)}(x)=\frac{N!}{d!}\left(\frac aN\right)^{N-d}
                         \left(1+\frac{ax}{N}\right)^d.
 \tag{1}
\]
All \(d\) roots are \(-N/a<0\). Thus the defining approximation
condition holds, with threshold \(N=d+1\). Moreover,
\[
 \gamma_0=c+1>0,\qquad \gamma_n=a^n>0\quad(n\ge1),
\]
so \(\phi\in\SLP\).

The coefficient criterion in \cite[Theorem 1.1]{W} gives a separate
check of membership: for every \(m\ge1\) and \(d\ge1\), the shifted
Jensen polynomial is
\[
 J_\gamma^{d,m}(X)=\sum_{k=0}^d\binom dk\gamma_{m+k}X^k
                 =a^m(1+aX)^d,
 \tag{2}
\]
whose roots are all negative.

For real \(x,y\), direct multiplication gives
\[
 |c+e^{a(x+iy)}|^2=c^2+e^{2ax}+2ce^{ax}\cos(ay).
\]
Consequently
\[
 L_j(\phi)(x)=\frac{2c(-1)^j a^{2j}}{(2j)!}e^{ax}
                         \qquad(j\ge1).                 \tag{3}
\]
For odd \(j\), every exponential Taylor coefficient of this function
equals \(-2ca^{2j+n}/(2j)!<0\). Eventual nonnegativity therefore
fails. In particular, \(L_1(1+e^x)=-e^x\notin\SLP\).
\end{proof}

\paragraph{Scope.}
The counterexample has strictly positive Taylor coefficients at every
index and satisfies both the direct approximation definition and the
shifted Jensen criterion. Repeated negative Jensen roots are allowed
in the source. The conclusion concerns the positive shifted class;
the negative exponential remains in the unsigned shifted class.

\paragraph{Task 2. Analyze polynomial perturbations.}
For \(a>0\) and real polynomials \(P\), characterize when
\(L_j(e^{ax}+P(x))\) belongs to \(\SLP\), and identify which
features of \(P\) control eventual coefficient signs.

\paragraph{Task 3. Investigate preservation after differentiation.}
Determine sufficient conditions on \(\phi\in\SLP\) under which,
for each \(j\), some \(N\) satisfies
\(L_j(\phi^{(n)})\in\SLP\) for every \(n\ge N\).

\paragraph{Task 4. Test strict shifted hyperbolicity.}
Determine whether counterexamples persist when every fixed-degree
shifted Jensen polynomial eventually has distinct negative roots.

\paragraph{Task 5. Examine the even operators.}
Decide whether the operators \(L_{2j}\), which preserve positivity
on the family in (3), preserve \(\SLP\) in general, or construct
counterexamples to that restricted assertion.

\begin{thebibliography}{1}
\bibitem{W} I.~Wagner,
\emph{On a new class of Laguerre--P\'olya type functions with applications
in number theory},
\href{https://arxiv.org/abs/2108.01827v2}{arXiv:2108.01827v2} (2022),
pp.~3--4 and Conjecture 4.3 on p.~11.
Published in \emph{Pacific Journal of Mathematics} \textbf{320} (2022),
177--192, \href{https://doi.org/10.2140/pjm.2022.320.177}
{doi:10.2140/pjm.2022.320.177}.
\end{thebibliography}
\end{document}
